% Omslagmotief biologie: DNA-dubbelhelix en een cel.
\foreach \x in {-6,-2,...,70} {
  \draw[white,opacity=.4,line width=.6pt] (\x,{22+13*sin(\x*7)}) -- (\x,{22-13*sin(\x*7)});}
\draw[kleur1!60!white,line width=1.6pt,domain=-8:72,samples=160] plot (\x,{22+13*sin(\x*7)});
\draw[kleur2!70!white,line width=1.6pt,domain=-8:72,samples=160] plot (\x,{22-13*sin(\x*7)});
\begin{scope}[shift={(96,20)}]
  \draw[kleur3!55!white,line width=1.5pt,rounded corners=8pt] (-20,-15) -- (18,-19) -- (22,12) -- (-8,22) -- (-24,6) -- cycle;
  \fill[kleur4!65!white,opacity=.8] (-2,2) circle (6);
  \fill[white] (0,3.5) circle (1.3);
  \draw[white,opacity=.6,line width=.8pt] (8,-9) ellipse (4 and 2);
  \draw[white,opacity=.6,line width=.8pt] (-14,-4) ellipse (3 and 1.6);
\end{scope}
